\chapter{The verification ladder}\label{ch:verification}

The model of the previous chapter is a set of equations and a set of choices, and neither is worth anything until it has been tested. A new kinetic code earns trust slowly. Dark-SPECTRAX was built as a ladder of tests, each rung resting on an exact solution, an independent kinetic dispersion relation or an independent code, and ordered so that a failure at one rung is diagnosed before the next is climbed. This chapter narrates the ladder: the algebraic rungs A00 to A10, the plasma benchmarks B00 to B12 and the numerical rungs C00 to C09. For each rung we give the reference, with the equation that defines it, and the measured agreement. The numbers come from the records indexed in \texttt{docs/results.md} and from the status table of the repository (\texttt{docs/status.md}); tests that are part of the pytest suite carry their identifier. Failed and amended gates are reported where they occurred, and a summary is collected in Table~\ref{tab:gatefail} at the end. The rule that governed all of it, stated in the working brief, is that a gate is declared in a committed script before the run, and if the gate turns out to be wrong the failure is recorded and a new follow-up gate is declared in a new commit before the next run. A gate is never edited after its result is seen.

\section{Algebra and exact solutions}\label{sec:algebra}

The A rungs involve no time integration of a nonlinear problem. They check that the equations of Chapter~\ref{ch:model} have been transcribed correctly, and they are in the pytest suite. A00 compares the Hermite moments and the Lorentz operator with an independent quadrature of the same integrals; the agreement is $10^{-12}$ or better, which also confirms the sign of the Levi-Civita term. A01 asks that with zero mixing the companion reproduce the parent: the right-hand sides agree to $10^{-14}$ and a short trajectory is identical. A02 checks the Proca vacuum: the dispersion relation $\omega^2=k^2c^2+\OmD^2$, the polarization and the two constraints. A03 integrates the cold homogeneous coupled oscillators of Eq.~\eqref{eq:coldroots} against a matrix exponential and against both roots of the quartic, to $10^{-11}$. A04 does the same at finite $k$, for the cold transverse and longitudinal branches taken from full determinants; the relative differences, $2\times10^{-5}$ and $10^{-4}$, are the size of the thermal correction at $v_t=0.005c$ and are reported as such. A05 is the static Yukawa limit, in which the dark scalar potential of a point charge is screened as $e^{-\OmD r/c}/r$. A06 verifies the symmetry under $\eta\to-\eta$. A07 tests the exact mean pump of Eq.~\eqref{eq:linlaw}, $\bar E=-\tfrac12E_0\omega t\sin\omega t$, to $10^{-9}$; this is exact in the discrete system only if the closure leaves indices $0$ to $2$ untouched, and the test would catch a closure that did not. A08 checks the eighth-order temporal convergence of fixed-step Dopri8 and the Hermite-order convergence of exact free streaming; the status table records that it uses exact solutions only, with no manufactured-source hook. A09 runs one to three populations on odd and even grids with ledger closure, and A10 tests the failure policy: a nonfinite state, an exhausted step budget, and a short stiff transient without a step floor each produce an explicit status instead of a silent wrong answer.

Two further tests belong to the algebraic group because they are exact. C00 draws random states that satisfy both Gauss laws and verifies, for particles, Maxwell and Proca separately, the energy-work theorems of Eq.~\eqref{eq:ledger}, the continuity equation, the preservation of both constraints and of $\BD=\nabla\times\AD$ by the right-hand side, and the vanishing of magnetic work. As a safeguard against tests that pass for the wrong reason, deliberate mutations of the companion (a flipped sign, a wrong source, a wrong mass potential, a wrong weight) were each caught by at least one test.

\section{Linear Landau damping}\label{sec:landau}

The first plasma benchmark, B00, is linear Landau damping of the Langmuir wave with and without mixing. The reference is the root of the longitudinal dispersion relation of Eq.~\eqref{eq:DL}, with the Maxwellian susceptibility written through the plasma dispersion function $Z$,
\begin{equation}
\chi_{\rm kin}(\omega,k)=\frac{1}{k^2\lambda_{De}^2}\Big[1+\zeta Z(\zeta)\Big],\qquad \zeta=\frac{\omega}{\sqrt2\,k\,v_{te}},
\label{eq:chikin}
\end{equation}
evaluated through the Faddeeva function of SciPy (\texttt{wofz}) with the causal continuation, and, for the dark case, the extra factor and the mixed branches of Eq.~\eqref{eq:DL}. The set-up has electrons with $v_{te}=0.1c$, a fixed neutralizing background, a density seed $10^{-4}$, the ordinary Gauss field plus the static Yukawa dark near field, and Dopri8 with $\mathrm{rtol}=10^{-10}$ and $\mathrm{atol}=10^{-16}$, without any closure. The complex frequency is fitted to $E_{x,k}(t)=\sum_jc_je^{-\ii\omega_jt}$ with the pair of Landau modes (plus the massive longitudinal branch for the dark runs) by nonlinear least squares.

\begin{table}[t]
\caption{B00: fitted complex frequency ($N_n=128$, $N_x=8$) against the independent kinetic root, ordinary and dark ($\eta=0.3$, $\OmD=\wpe$). Relative error of $\omega$, absolute error of the damping rate, spread over $N_n=64,128$ and $N_x=5,8$, and the fit window in $\wpe^{-1}$.}\label{tab:b00}
\small
\resizebox{\ifdim\width>\textwidth\textwidth\else\width\fi}{!}{%
\begin{tabular}{@{}clccccc@{}}\toprule
$k\lambda_{De}$ & model & fit & root & rel.\ error & $\Delta\gamma$ & window\\\midrule
0.2 & ord. & $1.063940-0.000127\ii$ & $1.063984-0.000055\ii$ & $8.0\times10^{-5}$ & $7.2\times10^{-5}$ & 8--34\\
0.2 & dark & $1.094251-0.000107\ii$ & $1.094261-0.000028\ii$ & $7.4\times10^{-5}$ & $8.0\times10^{-5}$ & 8--34\\
0.3 & ord. & $1.159878-0.012605\ii$ & $1.159846-0.012620\ii$ & $3.0\times10^{-5}$ & $1.5\times10^{-5}$ & 8--24\\
0.3 & dark & $1.191599-0.009619\ii$ & $1.191574-0.009642\ii$ & $2.8\times10^{-5}$ & $2.3\times10^{-5}$ & 8--24\\
0.5 & ord. & $1.415651-0.153354\ii$ & $1.415662-0.153359\ii$ & $8.6\times10^{-6}$ & $5.9\times10^{-6}$ & 8--14\\
0.5 & dark & $1.447808-0.139147\ii$ & $1.447817-0.139149\ii$ & $6.8\times10^{-6}$ & $1.9\times10^{-6}$ & 8--14\\
0.7 & ord. & $1.673852-0.392174\ii$ & $1.673866-0.392401\ii$ & $1.3\times10^{-4}$ & $2.3\times10^{-4}$ & 4--10\\
0.7 & dark & $1.707528-0.368745\ii$ & $1.707571-0.368955\ii$ & $1.2\times10^{-4}$ & $2.1\times10^{-4}$ & 4--10\\\bottomrule
\end{tabular}}
\end{table}

Table~\ref{tab:b00} shows agreement of the complex frequency to $1.3\times10^{-4}$ over the 32 runs of the study (worst case at $k\lambda_{De}=0.7$ where the damping is so strong that only a short window is available), the spread over the Hermite and spatial resolutions is at the level $10^{-9}$, and the ledger defect is $\le3\times10^{-18}$. The mixing moves the roots by $+2.7\,\%$ in frequency and $-23.6\,\%$ in damping at $k\lambda_{De}=0.3$ (the fitted rates give $-23.7\,\%$), and by $+2.3\,\%$ and $-9.3\,\%$ at $0.5$. The table also exposes the one place where the test is honest about its limit. At $k\lambda_{De}=0.2$ the damping rate is $5\times10^{-5}$ and changes the amplitude by about $0.1\,\%$ over the 26 time units that the Hermite front allows before it reaches index 64; the fitted rate is off by about $7\times10^{-5}$ absolute, a factor two, so only the frequency is validated there and the damping is recorded as unresolved. The window is limited by recurrence. In a free-streaming Maxwellian whose width matches the basis, the coefficients obey $|g_n|^2\propto[(kv_tt)^2]^ne^{-(kv_tt)^2}/n!$, a Poisson distribution in $n$ with mean $(kv_tt)^2$, so the front reaches the last retained index when $(kv_tt)^2\sim N$. This gives $t\sim23$ for $N=128$ at $k\lambda_{De}=0.5$, consistent with the window 8--14 chosen before the front reaches the edge, and the fit-window sensitivity listed in the record (relative errors of $10^{-3}$ for windows starting at $t=2$ that fall to $10^{-5}$ for windows starting at $t=6$) shows the ballistic transient, not the truncation, as the early bias.

\section{Streaming instabilities}\label{sec:streaming-inst}

The unstable counterparts are the two-stream and bump-on-tail instabilities, B02 and B03, for which the growth rate is the root of the same dispersion relation with drifting Maxwellians in the susceptibility. The electrostatic inputs are mapped with $v_t/c=0.1$, $\eta=0.3$, $\OmD=\wpe$, a seed $10^{-13}$, and the growth is the slope of $\ln|E_k|$ over a 25-unit window ending where $|E_k|$ reaches $10^{-5}$.

\begin{table}[t]
\caption{B02/B03: fitted growth rate and the kinetic root, ordinary and dark ($\eta=0.3$, $\OmD=\wpe$). The two-stream rows are unchanged between $N_n=32$ and $64$; the bump-on-tail rows are $N_n=128$.}\label{tab:b02}
\small
\resizebox{\ifdim\width>\textwidth\textwidth\else\width\fi}{!}{%
\begin{tabular}{@{}llccc@{}}\toprule
case & model & fit & root & difference\\\midrule
two-stream, $v_t=0.1$, $k=0.6$ & ord. & 0.3490944 & 0.3490944 & $-1.1\times10^{-8}$\\
two-stream, $v_t=0.1$, $k=0.6$ & dark & 0.3625230 & 0.3625237 & $-6.8\times10^{-7}$\\
two-stream, $v_t=0.3$, $k=0.4$ & ord. & 0.2616155 & 0.2616156 & $-2.7\times10^{-8}$\\
two-stream, $v_t=0.3$, $k=0.4$ & dark & 0.2664469 & 0.2664469 & $6.1\times10^{-9}$\\
bump-on-tail, $k=0.3$ & ord. & 0.1980980 & 0.1980980 & $1.6\times10^{-8}$\\
bump-on-tail, $k=0.3$ & dark & 0.2085951 & 0.2085951 & $1.6\times10^{-8}$\\\bottomrule
\end{tabular}}
\end{table}

The agreement in Table~\ref{tab:b02} is better than $10^{-6}$ throughout, and the bump-on-tail difference falls from about $7\times10^{-7}$ at $N_n=64$ to $2\times10^{-8}$ at $N_n=128$, so the residual is a Hermite-resolution error of the cold bump, not of the method. The mixing raises each growth rate, by $3.8\,\%$ for the first row, but this is one deliberately large coupling at one wavenumber per case and not a map of the growth band. With the benchmarks B00 and B02/B03 in hand the pin of the parent was refreshed several times; all values other than timings were bit-identical across the refreshes, with two ledger defects of $6\times10^{-11}$ moving by $10^{-20}$.

\section{Independent grid solvers}\label{sec:grid}

A test that shares every convention with the code under test proves little, so the ordinary benchmarks were repeated against a NumPy semi-Lagrangian Vlasov--Poisson solver, imported with provenance in \texttt{studies/refs}, and against Faddeeva roots. Table~\ref{tab:gridref} is the comparison: the Hermite fit, a Richardson-extrapolated grid fit, and the root. The grid and the Hermite fit agree to $10^{-9}$ at $k\lambda_{De}=0.5$, and to $2\times10^{-6}$ at $0.3$, where the damping difference ($-2.3\times10^{-6}$) slightly exceeds the stated grid uncertainty ($1.6\times10^{-6}$) while the Hermite--root difference is $1.1\times10^{-7}$. The two-stream growth rates agree with the grid to $4.4\times10^{-6}$ or better. For the dark sector the same idea required a new reference, a grid Vlasov--Ampère--Proca solver extending the NumPy one, which agrees with the Hermite echo amplitude to $2\times10^{-5}$ (Sec.~\ref{sec:nonlinear}). The independent reference code of the Gkeyll family enters only through imported $p=2$ records, since the Gkeyll documentation was not read.

\begin{table}[t]
\caption{Hermite against the independent semi-Lagrangian grid and the Faddeeva root (ordinary plasma). Differences are Hermite minus grid.}\label{tab:gridref}
\small
\resizebox{\ifdim\width>\textwidth\textwidth\else\width\fi}{!}{%
\begin{tabular}{@{}lcccc@{}}\toprule
case & Hermite fit & grid (Richardson) & root & Hermite $-$ grid\\\midrule
B00 $k\lambda_{De}=0.5$, $N=128$ & $1.4156617-0.1533591\ii$ & $1.4156617-0.1533591\ii$ & $1.4156619-0.1533595\ii$ & $1.3\times10^{-9}$, $1.6\times10^{-8}$\\
B00 $k\lambda_{De}=0.3$, $N=128$ & $1.1598463-0.0126203\ii$ & $1.1598444-0.0126180\ii$ & $1.1598465-0.0126204\ii$ & $1.9\times10^{-6}$, $-2.3\times10^{-6}$\\
B02 $v_t=0.1$, $k=0.6$ & 0.3490982 & 0.3491004 & 0.3490944 & $-2.1\times10^{-6}$\\
B02 $v_t=0.3$, $k=0.8$ & 0.2823913 & 0.2823949 & 0.2823968 & $-3.6\times10^{-6}$\\\bottomrule
\end{tabular}}
\end{table}

\section{Ion-acoustic waves and the Buneman instability}\label{sec:ia-buneman}

With mobile ions the electrostatic response contains a second branch. B04 measures ion-acoustic damping for $T_e/T_i=1,5,20$ at $k\lambda_{De}=0.1$, $0.2$, $0.3$ with $m_i/m_e=1836$, a seed $10^{-6}$ and the hypercollision $\nu=1$, by fitting the complex frequency on $E_k$ and comparing with the kinetic root with both species in the susceptibility. At $N_n=128$ the frequency error is $4$--$5\times10^{-3}$ at $T_e/T_i=1$ (worst $7.3\times10^{-3}$ for the dark runs), $1.3$--$1.4\times10^{-3}$ at 5 and $1.0\times10^{-3}$ at 20. The damping error is up to $6\times10^{-3}$ at $T_e/T_i=1$, $1.5\times10^{-2}$ at 5, and $8\,\%$ weak at 20. The last is not resolved: the damping there corresponds to 0.3 e-foldings over the fit window, and is not claimed. Going from $N_n=32$ to 64 removes the electron-recurrence error at $T_e/T_i=5$ (damping errors up to $0.2$ at 32). Twelve runs at $N_x=8$ and $T_e/T_i\ge5$ exhausted the step budget without a nonfinite field and are recorded as failures; the $N_x=8$ runs that did finish reproduce $N_x=5$ to $5\times10^{-8}$. Mixing shifts the roots by at most $0.25\,\%$ in frequency and $0.9\,\%$ in damping.

The Buneman instability, B05, is a drift of electrons at $u_0=0.1c$ through ions at rest, with $v_{te}=0.1u_0$ and $T_i=T_e$. Because a periodic box cannot sustain a net current, the $k=0$ current is cancelled by an inert return-current carrier of charge $+1$ and mass $10^6m_e$ drifting with the electrons, whose susceptibility is kept in the roots, and the residual mean field stays below $9\times10^{-7}$ for a seed of $10^{-10}$. The cold-fluid scaling $\gamma_{\max}=(\sqrt3/2)(m_e/2m_i)^{1/3}\wpe$ is $0.148$ for $m_i/m_e=100$ and $0.0561$ for 1836, and the kinetic roots at the two sampled wavenumbers are $66$--$95\,\%$ of it, consistent with but not a test of the fluid scaling. The fitted growth rates in Table~\ref{tab:b05} agree with the kinetic roots to $1.5\times10^{-4}$ in all 24 runs, and the residual is the fit window (sub-window spread $2\times10^{-6}$ to $1.5\times10^{-5}$), not the truncation: $N_n=32$ and $N_x=8$ change the rate by at most $2.2\times10^{-7}$. Mixing changes the growth by $-0.8$ and $+1.5\,\%$ at $m_i/m_e=100$ and by $-8.4$ and $-3.9\,\%$ at 1836, and lowers the real frequency by $12$ to $36\,\%$.

\begin{table}[t]
\caption{B05: Buneman growth rate against the kinetic root; $N_n=16$, $N_x=5$. Root columns give $\omega_r+\ii\gamma$.}\label{tab:b05}
\small
\resizebox{\ifdim\width>\textwidth\textwidth\else\width\fi}{!}{%
\begin{tabular}{@{}rcccc@{}}\toprule
$m_i/m_e$ & $ku_0$ & model & root & rel.\ error of $\gamma$\\\midrule
100 & 0.9 & ord. & $0.056353+0.121750\ii$ & $-5.9\times10^{-5}$\\
100 & 1.0 & dark & $0.076413+0.134947\ii$ & $-1.6\times10^{-6}$\\
1836 & 0.9 & ord. & $0.008270+0.040672\ii$ & $-1.0\times10^{-4}$\\
1836 & 1.0 & dark & $0.018058+0.051220\ii$ & $+1.5\times10^{-4}$\\\bottomrule
\end{tabular}}
\end{table}

\section{Nonlinear Landau damping and the echo}\label{sec:nonlinear}

The first nonlinear rung, B01, is a Langmuir wave of finite amplitude, $k=0.3$ and $\varepsilon=0.05$, in the electrostatic units of the benchmark (with $\beta=0.1$). O'Neil~\cite{ONeil1965} showed that a wave of fixed amplitude traps resonant electrons, so that linear theory fails after the inverse deep-trapped bounce frequency $\tau=(kE)^{-1/2}$ in these units, and that the damping rate oscillates and phase-mixes to zero thereafter. Our amplitude $E=\varepsilon/k=0.1667$ gives $\tau=4.47$, $2\pi\tau=28.1$ and $\gamma_L\tau=0.056$, so the constant-amplitude assumption holds at onset. The envelope of the grid solution has a first minimum at $t=31.45$ and a maximum at 65.6, separated by 34.2, and $2\pi\tau$ evaluated with the measured envelope at those times is 39.7 and 33.1; the rebound timescale is therefore the O'Neil bounce scale. O'Neil gives no prediction for the extremum times of a decaying wave, so B01 remains a grid-against-Hermite comparison and is not promoted to a theory benchmark. At $N=1024$ with $\nu=0$ the Hermite run reproduces the grid's first minimum at $t=31.45$ and stays within $10\,\%$ of the grid envelope to $t=100$, while $N=512$ departs at $t=85$.

A correction belongs here, because a wrong statement made it into an early record. The imported reference stated that a hypercollision $\nu=1$ at $N=256$ and 512 removes the trapping rebound. This rested on a prominence-based extremum detector that also finds no extremum in the $\nu=0$, $N=512$ run. A detector-free re-measurement of the envelope ratio $R=\ln[\mathrm{env}(65.6)/\mathrm{env}(31.45)]$ at the grid extremum times, relative to the grid value $0.321$, gives $R/R_{\rm grid}=0.82$ and $0.93$ for $\nu=1$ at $N=256$ and 512, $0.88$ for $\nu=0$ at 512 and $1.13$ at 1024; and for $N=256$, $\nu=0$ it is $-0.38$, because recurrence corrupts the curve. The statement was retracted. The hypercollision keeps the rebound; what damages the late evolution at low $N$ is recurrence.

The echo, B06, tests the nonlinear phase-space structure. Two impulsive perturbations at wavenumbers $k_1$ and $k_2$, applied at $t=0$ and $t=\tau$, leave a second-order term $\propto\exp[\ii(k_2-k_1)x+\ii k_2v\tau-\ii(k_2-k_1)vt]$ whose velocity dependence vanishes at $t=\tau k_2/(k_2-k_1)$, so that an echo appears at $k_3=k_2-k_1$ (Gould, O'Neil and Malmberg~\cite{Gould1967}). For $k_1=1$, $k_2=1.5$ and $\tau=10$ the echo is at $t=30$ with $k_3=0.5$. The grid solution has its echo at $t=29.05$, 3\,\% early, which is the finite-damping shift that the ballistic estimate omits; the Hermite echo at $N=512$ is at 29.05 with an amplitude $+0.002\,\%$ of the grid value, $1.15875\times10^{-3}$ against $1.15873\times10^{-3}$. The kick is an external velocity shift applied at a restart, exact in the truncated basis, and its kinetic energy ($6.25\times10^{-6}$) is an external term outside the ledger. At $N=256$ the echo peak is reproduced but the recurrence at $k_1$ contaminates the curve by $8\,\%$. With the hypercollision $\nu=1$ the echo amplitude is lowered by $3.2\,\%$, so the closure is not free of consequences for fine phase-space structure. The dark runs show a first envelope minimum at $t=30.5$ in B01, earlier than the ordinary one, and an echo amplitude lowered by $3.3\,\%$ by the mixing. The Hermite and grid echo amplitudes agree to $2.0\times10^{-5}$ for the ordinary plasma and, with the dark grid reference added later, $1.6\times10^{-5}$ for the dark one.

\section{Magnetized plasmas and the Proca flux}\label{sec:magnetized}

B08 tests the magnetic force. The first part is the rotation in a uniform field. A drifting anisotropic Maxwellian ($\mathbf U_0=[0.02,-0.01,0.015]$, thermal speeds $[0.01,0.02,0.015]$) in a field of magnitude $0.5$ along $z$, $x$, $(1,1,1)$ and $(0.3,-0.5,0.8)$, for both signs of charge and two densities, must gyrate exactly: the mean flow, the mean field and the full covariance follow $P(t)=RP_0R^T$ with $R$ the rotation, and the mean-flow/mean-field subsystem has an exact six-by-six matrix exponential. In 16 runs the worst relative errors are $5.2\times10^{-12}$ for the mean flow, $4.5\times10^{-12}$ for the field, $2.1\times10^{-11}$ for the covariance and $2.2\times10^{-11}$ for the gyration frequency at density 1, and the kinetic energy changes by up to $49\,\%$ entirely through the electric work, so $K-K_0-W_{\rm em}$ closes to round-off and the magnetic work is zero. Two of the gates failed as declared at the lower density $10^{-4}$: the gyration frequency (gate $10^{-8}$, worst $1.2\times10^{-7}$ in 8 of 16 runs) and the ledger (gate $10^{-10}$, worst $3.3\times10^{-10}$ in 8 of 16), and a third was mis-specified; the causes recorded are the fit window, a small-$K_0$ ratio and a wrong $K_\perp$ definition. They are documented, not amended.

The second part compares waves. Electrons with $\wpe=1$ and $|\Omega_{ce}|=0.5$, $k$ parallel and perpendicular to the field at $k=0.5$ and $1.1$, ordinary and dark ($\eta=0.3$, $\OmD=1.2\wpe$): the reference is the set of eigenvalues of the full cold field-current operator, with every eigenfrequency also annihilating the conductivity-based field system (smallest singular value over largest $\le3\times10^{-15}$). Each positive-frequency eigenvector seeds a Hermite run of amplitude $10^{-5}$ that is fitted with a single mode. All 120 runs succeed and all 40 cold R, L, whistler, Langmuir, O and X eigenmodes, ordinary and dark, are recovered within $10^{-3}$ at $v_t=0.01$ (worst $9.9\times10^{-4}$), with a refinement in $N_n$ changing at most $10^{-6}$. The remaining error is the warm correction: halving $v_t$ divides it by $4.0$, the $v_t^2$ scaling. Near $k=1.1$ the R and upper-X branches cross $\sqrt{k^2+\OmD^2}$ and their eigenvectors are strongly mixed. The kinetic ledger closes to $4\times10^{-10}$ of the wave energy and the Gauss residual to $10^{-13}$. The third part is the Proca flux: packets of five Proca modes (transverse, longitudinal, mixed) in a cold plasma with a background field are evolved to $t=5$ and the pointwise balance $\partial_tu_D+\partial_xS_D+P_D=0$ is checked on 256 points. The residual relative to $\max|\partial_tu_D|$ is $10^{-14}$, $3\times10^{-14}$ and $2\times10^{-14}$ for the three packets, against $0.02$, $1.15$ and $1.37$ if the term $\OmD^2\phiD A_{D,x}$ is dropped from the flux. A warm (Bernstein) reference for the lower X branch near $2|\Omega_{ce}|$ and a nonuniform background field are open.

Weibel and filamentation instabilities, B09, test the dark coupling in the transverse sector. In 1D3V a bi-Maxwellian plasma with temperature anisotropy, and a counter-streaming pair for the current filamentation, drive a magnetic perturbation at rate $\gamma(k)$ given by an independent kinetic transverse dispersion relation with the Proca coupling included. Table~\ref{tab:b09} lists the ordinary Weibel cases. The base level $(N_x,N_n,N_m)=(5,32,4)$ reproduces the rate to better than $10^{-6}$ at $k=0.4$ and $0.8$, and the error grows at higher $k$ because the Hermite order needed rises: at $k=1.6$ one needs $N_n=128$, at $2.0$ also 128 to 256, and at $k=2.4$ the declared refinement step $N_n=256\to512$ changes the rate by $1.4\times10^{-3}$ and fails the gate, although the finest result is $6.2\times10^{-5}$ from the root. The dark runs behave alike for $\OmD=0.3$ (near the transverse Proca branch) and $1.2$, with the exception of one spurious fit-mode pick at $k=0.4$, $N_x=9$ (error $0.64$), a failure of the fitting procedure and not of the code, recorded as such. The ledger is $\le1.1\times10^{-7}$. The part-one overall gate and the ladder gate are therefore reported as failed.

\begin{table}[t]
\caption{B09, linear Weibel instability, ordinary plasma: relative error of the fitted growth rate $\gamma$ against the kinetic root at the base level, the velocity-refined level, and the $N_n$ ladder. Gate: the finest refinement step is below $10^{-4}$.}\label{tab:b09}
\small
\resizebox{\ifdim\width>\textwidth\textwidth\else\width\fi}{!}{%
\begin{tabular}{@{}lccccccc@{}}\toprule
$k$ & $\gamma_{\rm root}$ & base & $v$-refined & $N_n=128$ & 256 & 512 & gate\\\midrule
0.4 & 0.065020 & $2.4\times10^{-7}$ & $2.4\times10^{-7}$ & & & & pass\\
1.2 & 0.107751 & $1.1\times10^{-5}$ & $5.0\times10^{-8}$ & & & & pass\\
1.6 & 0.095873 & $6.0\times10^{-4}$ & $1.0\times10^{-5}$ & $6.8\times10^{-8}$ & $6.8\times10^{-8}$ & $6.8\times10^{-8}$ & pass\\
2.0 & 0.072070 & $1.4\times10^{-2}$ & $1.3\times10^{-3}$ & $4.6\times10^{-5}$ & $1.3\times10^{-7}$ & $1.3\times10^{-7}$ & pass\\
2.4 & 0.040307 & $1.5\times10^{-1}$ & $5.5\times10^{-2}$ & $1.3\times10^{-2}$ & $1.5\times10^{-3}$ & $6.2\times10^{-5}$ & fail\\\bottomrule
\end{tabular}}
\end{table}

The nonlinear Weibel runs carry the caveat that matters. A bi-Maxwellian seeded with $B_z=10^{-6}$ at $k_0=1.2$ grows at $\gamma=0.1078$ and should saturate when the magnetic trapping frequency reaches the growth rate, at $B_{\rm sat}=\gamma^2/(k_0v_\perp)=4.56\times10^{-2}$. With no closure the run fails in the chunk ending $t=112$, with the top third of the Hermite ladder holding more than a tenth of the energy by $t=110$. The frozen-field spectrum at $t=100$ has a growing mode with real part $1.31$, twelve times the Weibel rate, all of it in the top third of the ladder: the truncation instability of Chapter~\ref{ch:limits}, here driven by $\vv\times\BB$ rather than $E$. The electric closure of Eq.~\eqref{eq:closure} barely changes it (growth $1.23$) because $E_x$ is small in Weibel; the magnetic closure removes it at $t=100$ (growth $0.0022$) and reduces it by a factor of ten at $t=108$. The magnetic closure at $c=1$ only postpones the collapse to $t\simeq148$. Declared criterion 1 and the gate G3a failed (the velocity-doubled run at $c=1$ fails before the first magnetic maximum). The peak result was obtained after post-hoc amendments, committed before their runs but after the failure was seen: under them the first maximum of the magnetic energy is at $t=115.5$--$120.5$ with $B_{\rm peak}=1.93$--$1.97\,B_{\rm sat}$ for the ordinary and the two dark models, unchanged to $0.5\,\%$ by doubling $N_x$, by $c=0.5$ and 2, and (at $c=2$) by doubling the velocity resolution. The ratio of the dark to the ordinary peak follows the bounce estimate to $2.3\,\%$ ($\OmD=0.3$) and $1.3\,\%$ ($\OmD=1.2$). Nothing is claimed after the first maximum, and filamentation was not run in the nonlinear regime.

\section{Langmuir modulational instability}\label{sec:b10}

The Langmuir pump of finite amplitude is unstable to sidebands, the modulational instability described by the Zakharov equations. Zakharov~\cite{Zakharov1972} writes the slowly varying envelope $E$ of the fast field, coupled to the ion density perturbation $\delta n$, as $\ii E_t+\tfrac32\wpe r_D^2E_{xx}=\tfrac{\wpe}{2n_0}\delta n\,E$ and $(\partial_t^2-c_s^2\partial_x^2)\delta n=\tfrac{1}{16\pi M}\partial_x^2|E|^2$ in Gaussian units for the one-dimensional case. For a pump at $k_0=0$ the static limit gives a growth rate $\gamma=\wpe[\tfrac34q\kappa^2r_D^2w-\tfrac94\kappa^4r_D^4]^{1/2}$ for sideband wavenumber $\kappa$, with $w=W/nT_e$ the pump energy ratio and $q=T_e/(T_e+T_i)$. Reading the paper in detail revealed a misprint: maximizing this expression gives $(\kappa r_D)^2=qw/6$ and $\gamma_{\max}=\tfrac14q\wpe w$, whereas the text prints $\tfrac1{16}q\wpe w$ and $(\kappa r_D)^2=w/6$. B10 therefore uses the unmaximized expressions directly (Eqs.~2.8 and 2.9 of the paper), and the code of the reference reproduces the maximum $\tfrac14q\wpe w$ for the static limit and $(\mu w/3)^{1/2}$ for the hydrodynamic one. Two references are built: R0, the Zakharov envelope growth, and R1, a semi-kinetic one with the exact Langmuir dispersion relation and kinetic ions, extended to the dark case by replacing the force by the field $E_f=gE$ with $g=1+\eta^2Q/(Q-\OmD^2)$ from Eq.~\eqref{eq:EDratio}, which modifies the sideband operator and the ponderomotive coupling.

The simulation uses a pump that is an eigenvector of the model itself ($w=0.04$, $k_0r_D=0.005$, $T_e/T_i=10$, $v_{te}=0.1c$), with one sideband pair per run seeded at relative amplitude $3\times10^{-4}$, hypercollision $\nu=100$ and $N_n=32$; for the dark model the pump is one of the two mixed Langmuir--Proca modes. For the ordinary plasma the fitted sideband growth is $-2.5$, $-1.8$ and $-3.8\,\%$ from R1 at $\kappa r_D=0.0125$, $0.02$, $0.04$ for $m_i/m_e=1836$ and $-6.9$, $-5.6$, $-5.7\,\%$ for $m_i/m_e=100$, inside the declared $15\,\%$ and $30\,\%$ against R0; R0 and R1 differ by less than $1\,\%$, so the fit tests the Zakharov structure itself. The pump changes by at most $0.25\,\%$ over the window, the ledger is $\le10^{-9}$, and the growth is $0.7$--$1.8\times10^{-3}\wpe$, a window of $2.3$ e-foldings.

The summary of B10 has more qualifications than the headline numbers suggest. The hypercollision is the weak point: $\nu=1,3,10,40$ all blew up in shakedown runs and $\nu=100$ is the smallest stable value found; with $\nu=300$ the ordinary growth changes by $-5.3\,\%$ ($m_i/m_e=1836$) and $-8.7\,\%$ (100), so the declared closure-sensitivity gate G2 (3\,\%) failed, and the ordinary agreement above carries a closure uncertainty of $5$--$9\,\%$. The $N_n=64$ refinements blew up, so no $N_n$-converged result exists. The dark upper mode agrees with R1 to $-9.7\,\%$ at the R1 maximum for $m_i/m_e=1836$ but is over-predicted by $16$--$40\,\%$ elsewhere (gate G1 passes at 10 of 21 dark points overall and at only one of five for $m_i/m_e=100$); the two points $j=7,8$ at $m_i/m_e=1836$ failed at $\nu=100$ with growth three to four times R1 followed by an explosive stage, an artifact of the closure since $\nu=300$ removes it, but no gated agreement exists there. The dark lower mode is predicted to grow so slowly that the windows cannot resolve it. An exact tangent-flow Floquet cross-check agrees with the fit to $-1.2\,\%$ for the ordinary case ($2.51\times10^{-3}$ against $2.48\times10^{-3}$) and is inconsistent with both fit and R1 for the dark case; the Floquet runs at $m_i/m_e=1836$ returned numbers contradicting stable simulations and are a failure of the fixed-step tangent solve. The qualitative statements that the dark modes move the unstable band to smaller $\kappa$ and reduce the rate at large $\kappa$ are supported by the measured ordering ($\gamma_{\rm dark}/\gamma_{\rm ord}=1.10$ against $1.19$ predicted at $j=5$; $0.58$ against $0.90$ at $j=10$ for $m_i/m_e=100$), but the dark envelope is incomplete away from its peak.

\section{Collisions}\label{sec:b11}

B11 verifies the Dougherty operator of Sec.~\ref{sec:collisions} against exact relaxation rates, in homogeneous uncharged cases where the collision operator is the whole dynamics (Dopri8 with $\mathrm{rtol}=\mathrm{atol}=10^{-12}$). The references are the exact rates of the implemented operator, derived in the script docstring from Eq.~\eqref{eq:lbo} and the moment equations. A pair of beams in one dimension with width mismatched to the basis must relax each mode as $C_n(t)=C_n(0)e^{-n\nu t}$ in the shifted-centred basis; three-dimensional anisotropy $T_i-\theta$ decays as $e^{-2\nu t}$; two species at different temperatures obey $\Delta T\propto\exp[-2(\nu_{sr}m_s+\nu_{rs}m_r)t/(m_s+m_r)]$; and a relative drift decays as $\exp[-(\nu_{sr}+\nu_{rs})t/2]$. Table~\ref{tab:b11} shows agreement between $4\times10^{-11}$ and $10^{-15}$ with the conserved moments held to $10^{-15}$. The non-Maxwellian part of the two-beam case falls from $0.191$ to $2.7\times10^{-10}$ by $t=5$. The relative mode error grows as $C_n$ decays toward the absolute tolerance, and the gate was restated from a whole-window relative error to an absolute plus windowed relative error after an uncommitted pre-run, which the script discloses. The cross-species coefficients of parent \#71 equal the LBO-EM form of Francisquez et al.~\cite{Francisquez2022} and not Greene's $\beta=0$ form, as discovered by comparing equations; LBO-G was then added as parent \#73.

\begin{table}[t]
\caption{B11: Dougherty relaxation against exact rates (homogeneous, uncharged).}\label{tab:b11}
\small
\resizebox{\ifdim\width>\textwidth\textwidth\else\width\fi}{!}{%
\begin{tabular}{@{}L{5.7cm}L{3.2cm}cc@{}}\toprule
case & reference & error & conservation\\\midrule
two beams, 1V, $N_n=96$, $\nu=1$ & $C_n(0)e^{-n\nu t}$ & $4.3\times10^{-11}$ (abs.) & $2\times10^{-17}$\\
anisotropy $T=(0.25,0.16,0.09)$, $\nu=0.5$ & $e^{-2\nu t}$ & $4.4\times10^{-11}$ & $5\times10^{-16}$\\
$T_s/T_r=5$, $m_r/m_s=4$ & exponential, see text & $1.5\times10^{-13}$ & $9\times10^{-16}$\\
same, $m_r/m_s=100$ & same & $1.7\times10^{-15}$ & $2\times10^{-15}$\\
drift $U=(0.3,-0.075)$, $m_r/m_s=4$ & $e^{-(\nu_{sr}+\nu_{rs})t/2}$ & $1.3\times10^{-13}$ & $4\times10^{-16}$\\\bottomrule
\end{tabular}}
\end{table}

Collisional Landau damping extends B00 with self collisions at $\nu=0.01$, $0.05$, $0.2$ for $k\lambda_{De}=0.3$ and $0.5$, ordinary and dark. The reference is the linear Dougherty--Vlasov root, derived in two independent ways: the Hermite tridiagonal with collision diagonal $\nu n$ for $n\ge3$ (the restoring terms cancel rows 1 and 2), solved at $N=4000$ and 8000, and a continuum form that Fourier transforms in $v$ and uses the Lenard--Bernstein kernel along $s\in[0,k/\nu]$ with the momentum and temperature restoring terms closed by a three-by-three system. The two references agree to $4\times10^{-12}$ in the root. The fitted frequency agrees within $5.7\times10^{-4}$ in every case (a typical value is $2.4\times10^{-5}$), with $N_n=128$ against 256 to $6\times10^{-7}$ and a ledger of $10^{-17}$. At $k\lambda_{De}=0.5$ the damping falls with $\nu$ (ordinary $0.1534$ at $\nu=0$ to $0.0823$ at $\nu=0.2$; dark $0.1391$ to $0.0762$), as the damping of the $n\ge3$ tail weakens the phase mixing. The largest fit errors are not from the reference, since the two references agree, and the short fit windows are the likely cause, untested. Two declared gates failed, both mis-specified: G0 compared the collisional root at $\nu=10^{-6}$ with the collisionless one to $10^{-6}$ when the physical change is $\mathcal O(\nu)$, and G3 required a per-state row ratio that is one by construction at $t=0$. The follow-up gates, declared in a new script, show the change of the root linear in $\nu$ (ratios $0.1001$, $0.1000$) and a conservation level of $3.3\times10^{-12}$ which is round-off of the $\mathcal O(350)$ $C_0$ row rather than a defect; the ledger itself closes to $10^{-17}$.

The remaining collisional rung, B12, forced oscillation with drag, is only partly done: the response $\dot J=\omega_L^2(\bar E+E_0\cos\omega t)-\nu_{\rm eff}J$ with $\nu_{\rm eff}=\tfrac{\nu_{ei}}{2}(1+m_e/m_i)$ and the absorbed power $P_{\rm ss}=\omega_L^2\nu_{\rm eff}\omega^2E_0^2/\{2[(\omega_L^2-\omega^2)^2+\nu_{\rm eff}^2\omega^2]\}$ were derived, and the 192-run absorption study built on them is reported in the paper (the kinetic power equals the linear Drude law, with a quadratic exponent in the drive $1.99\pm0.10$ and the steady state within $0.4\,\%$ at $\nu_{ei}=3\times10^{-2}$). The collision rate is constant, so heating does not change the rate, and a cosmological statement would need the Coulomb scaling $T_e^{-3/2}$.

\section{Numerical rungs: conservation, recurrence, PIC, restart, GPU and gradients}\label{sec:crungs}

The C rungs examine the numerics. C00 was covered with the algebra. C01, the long-run phase error, is partial: the midpoint error $-\omega^3\Delta t^2t/12$ was derived and Dopri8's order was tested in A08, but no long-run measurement was made. C02 is refinement, done for the strong-drive set-up from $N_x=8$ to 128 and $N_n=32$ to 256 and presented in Chapter~\ref{ch:limits}.

C03 is recurrence and closure, which was established as a finding, not as a number: the Hermite coefficient front of free streaming reaches the last retained index at $(kv_tt)^2\sim N$ (Sec.~\ref{sec:landau}), a signal returns as a recurrence, and the closure of Eq.~\eqref{eq:closure} suppresses it at a price in the echo (Sec.~\ref{sec:nonlinear}).\label{sec:recurrence} C04, the remap, is covered by the parent tests of \#56--\#60 and used in the pump-frame lanes.

C05 compares Dark-SPECTRAX with the PIC code Dark-JAX-in-Cell~\cite{DarkJIC} on a matched Landau case: $\sigma=0.05c$, $k\lambda_{De}=0.5$, $\eta=0.3$, $\OmD=kc$, density seed $0.01$, Yukawa-consistent start, and the same fit rule (neighbouring maxima of $|E_k|$ on $2\le t\le12$). The two codes share the field conventions and the dispersion relation, whose roots agree to $4\times10^{-14}$, but not the kinetic discretization, the field solver, the integrator or the loading. Table~\ref{tab:c05} shows the PIC fits approaching the Hermite fit under joint refinement of cells and particles: the RMS log-amplitude difference falls from $0.275$ at 32 cells to $0.115$ at 128 cells, but the 128-cell fits remain $1$--$2\,\%$ from the Hermite fits in damping. The maxima rule with a $1\,\%$ seed is itself biased by about $1\,\%$ from the linear root, so the comparison is between like fits. This is not a claim of PIC convergence. The PIC noise limits the comparison to about $1\,\%$, which is an observation about the PIC, not about the Hermite solver whose result stands within $10^{-4}$ of the root (Table~\ref{tab:b00}).

\begin{table}[t]
\caption{C05: Landau fits (maxima rule) of the Hermite solver and of the PIC at increasing resolution, ordinary and dark; the last column is the RMS difference of $\ln|E_k|$ with Hermite.}\label{tab:c05}
\small
\resizebox{\ifdim\width>\textwidth\textwidth\else\width\fi}{!}{%
\begin{tabular}{@{}lccc@{}}\toprule
run & ordinary & dark & RMS diff (ord./dark)\\\midrule
Hermite $N_n=128$ & $1.41354-0.15497\ii$ & $1.42962-0.14745\ii$ & ---\\
PIC 32 cells, 40000 particles & $1.38528-0.17775\ii$ & $1.41280-0.17364\ii$ & 0.275 / 0.222\\
PIC 64 cells, 80000 particles & $1.39891-0.15525\ii$ & $1.41671-0.15038\ii$ & 0.219 / 0.203\\
PIC 128 cells, 160000 particles & $1.41008-0.15175\ii$ & $1.42499-0.14522\ii$ & 0.115 / 0.131\\
root & $1.41566-0.15336\ii$ & $1.43243-0.14580\ii$ & \\\bottomrule
\end{tabular}}
\end{table}

C06 asks for a grid or Gkeyll comparison. It is done with our own NumPy semi-Lagrangian grid solvers, including the Vlasov--Ampère--Proca extension, and partial for Gkeyll, which enters only through imported records. C07 is restart: a chunked run restarted from a saved state equals the single run with a continuous ledger. A caution came from the later collision study. The nonlinear $\nu=0$ rerun at the new parent differed from the older record by $6\times10^{-6}$ in $\Delta K_e$ against a gate of $10^{-6}$; the cause, found in a follow-up declared after the failure, is that the restarts of the adaptive integrator reset the PID controller and the first step, and the older record used a chunk of 240 time units against 200 for the rerun. With the chunk matched the rerun reproduces the record exactly and with chunks 200 and 120 the deviations exceed $10^{-8}$, as the follow-up gate declared. The gate failure is recorded, and the parent change is exonerated.

C08 compares a GPU with a CPU. On an RTX A4000 in float64, the H05 Hermite case $\vq=0.1$, $N_x=32$, $N_n=64$, $c=4$, $T=100$ gives a maximum relative difference of $9.7\times10^{-9}$ over $\Delta K_e$ ($9.2\times10^{-9}$), $\Wext$ ($5.5\times10^{-10}$), $T_x$ ($4.6\times10^{-11}$) and $|E_k|^2$ for $k\le8$ ($9.7\times10^{-9}$), against a declared gate of $10^{-6}$, with 476 steps on CPU and 478 on GPU. The two trajectories are adaptive and need not be identical, so this is agreement of two valid solutions, not bitwise identity. The GPU integration took 1.0 s against 3.9 s for four CPU threads. The comparison covers one short case, and nothing else in the physics was compared across devices.

C09 differentiates. A solver written in JAX can be differentiated end to end, and the question is whether gradients of converged dark observables are correct. The objectives are the two fitted frequencies of the cold coupled oscillator, the accumulated dark work $W_D(T=60)$, and $|E_k(10)|^2$ of a dark Landau run, with respect to $\eta$, $\OmD$, $n_0$, $T$ and the seed amplitude (19 derivatives). The references are an analytic derivative (the quartic roots, or the matrix exponential), forward mode, and centred differences over $h/|p|=10^{-1}$ to $10^{-7}$. Reverse mode equals forward mode to $2\times10^{-13}$ at worst, agrees with the analytic derivative to $5\times10^{-10}$ or better, and the finite-difference plateau reaches $6\times10^{-9}$ or better; the finite-difference ladder shows the textbook $h^2$ truncation followed by round-off. Table~\ref{tab:c09} gives a sample. The cost of reverse mode is 3.9 to 6.2 times that of the forward solve.

\begin{table}[t]
\caption{C09: reverse-mode derivatives against forward mode, analytic references and the best finite difference.}\label{tab:c09}
\small
\resizebox{\ifdim\width>\textwidth\textwidth\else\width\fi}{!}{%
\begin{tabular}{@{}llcccc@{}}\toprule
objective & parameter & reverse AD & fwd vs rev & best FD & AD vs analytic\\\midrule
$\omega_1$ & $\eta$ & $-3.3397956\times10^{-1}$ & $9\times10^{-15}$ & $10^{-11}$ & $8\times10^{-15}$\\
$\omega_2$ & $\Omega_D$ & $8.0915015\times10^{-1}$ & $5\times10^{-16}$ & $2\times10^{-10}$ & $10^{-15}$\\
$W_D(60)$, adaptive & $\eta$ & $-1.6419597\times10^{-4}$ & $6\times10^{-15}$ & $6\times10^{-9}$ & $2\times10^{-11}$\\
$|E_k(10)|^2$ & $\eta$ & $-2.2474371\times10^{-12}$ & $4\times10^{-15}$ & $2\times10^{-10}$ & ---\\
$|E_k(10)|^2$ & $T$ & $-3.2365028\times10^{-10}$ & $10^{-15}$ & $3\times10^{-12}$ & ---\\\bottomrule
\end{tabular}}
\end{table}

The study also found a defect upstream. With a constant step size and many \texttt{SaveAt(ts)} times that coincide with step ends, Diffrax's \texttt{RecursiveCheckpointAdjoint} returns wrong gradients, by a few percent, and the wrong value depends on the number of checkpoints and on \texttt{max\_steps}. The minimal case is the oscillator $y''=-ay$ with loss $\sum_iy(t_i)\cos t_i$ over 601 saves at $\mathrm{d}t=0.01$: the exact gradient is $331.6732904$, \texttt{DirectAdjoint} and \texttt{ForwardMode} reproduce it, and the default checkpointing returns $324.378$ ($2.2\,\%$ off). Changing \texttt{max\_steps} from 620 to 640, 700, 1000, 4096 changes the error to $2.9\times10^{-4}$, $1.4\times10^{-2}$, $1.1\times10^{-2}$, and then exact; checkpoint counts 2, 4, 8, 30, 300, 600 are exact, and 16, 20, 25, 34, 35, 40, 50, 100 are wrong; shifting the save times by $10^{-12}$, by $10^{-9}$, or by $\mathrm{d}t/2$ restores exactness. Final-state objectives with two saves and a short run of 11 saves are exact. The cause was not isolated. The case was reported upstream as patrick-kidger/diffrax\#788, and a regression test in SPECTRAX \#68 documents the failure with a strict expected failure. Our production gradients therefore use \texttt{DirectAdjoint} or forward mode. Gradients through the pump-frame remap and \texttt{run\_adaptive}, through a fitted damping rate, and through nonlinear-regime objectives are not covered.

\section{Beyond the ladder}\label{sec:beyond}

Two groups of rungs sit outside the A, B and C lists and are summarized in the paper and the status table rather than here. The H rungs reproduce the setting of Hook, Huang and Shalaby: H00, the homogeneous resonant mean field, agrees with the exact oscillator to $1.1$--$1.4\times10^{-9}$ in $\bar E$ and $1.6\times10^{-11}$ to $3.3\times10^{-11}$ in $\Wext$ up to $\wpe t=1000$, even at quiver speeds of $15\,\vte$; the finite-$k$ pilots H01 and H02 satisfy the mean-pump identity to $10^{-13}$; the finite reservoir H07 follows a $4\times4$ matrix exponential to $10^{-9}$. The conversion at a density ramp, R3, uses a cold-fluid Maxwell--Proca time-domain reference (energy drift $3.5\times10^{-10}$, refinement $4.5\times10^{-5}$, flux budget $3.4\times10^{-6}$ for transverse cases) against which the kinetic 1D3V packets crossing a tanh ramp through the resonance $\wpe=\OmD$ agree to $4\times10^{-4}$ in 16 transverse runs with a ledger $\le2\times10^{-10}$ of the incident energy; the simple Landau--Zener law tested against them failed its declared check in two of four gated cases, with an exponent too large by a factor traced to the ramp parameter. These are reported with the physics results.

\section{What failed on the ladder}\label{sec:failed-ladder}

The ladder is not a record of unbroken success, and Table~\ref{tab:gatefail} collects the failures and amendments of this chapter in one place. They fall into three kinds. A few are real limits of the model or of its closure (B09 at $k=2.4$, B10 closure sensitivity, B04 damping at $T_e/T_i=20$, B01 at low resolution). Several are gates that were written wrongly and were replaced by follow-ups declared before their runs (B11 G0 and G3, B08 gates on $K_\perp$, the relativistic spectrum c). One is a retraction of a stated result (the B01 hypercollision statement). The ladder is credible only if all of them are listed with their evidence, and none has been edited after the fact.

\begin{table}[t]
\caption{Failed, amended and retracted items on the verification ladder.}\label{tab:gatefail}
\small
\resizebox{\ifdim\width>\textwidth\textwidth\else\width\fi}{!}{%
\begin{tabular}{@{}lL{6.4cm}L{5.0cm}@{}}\toprule
Rung & Failure & Disposition\\\midrule
B00 & damping at $k\lambda_{De}=0.2$ unresolved inside the Hermite window & only the frequency validated\\
B01 & prominence detector missed minima; early claim that $\nu=1$ removes the rebound & retracted; detector-free remeasurement\\
B04 & damping $8\,\%$ weak at $T_e/T_i=20$; $N_x=8$ runs exhausted the step budget & recorded, not claimed\\
B08 & rotation gates at density $10^{-4}$ (frequency, ledger, $K_\perp$) & fit-window and definition causes documented\\
B09 & $k=2.4$ step gate; one spurious fit pick; nonlinear criterion 1 and G3a & peak under post-hoc amendments 1 and 3 only\\
B10 & closure sensitivity G2 ($5$--$9\,\%$); $N_n=64$ blow-up; dark G1 11 of 21 fail & dark agreement only at the peak\\
B11 & G0 and G3 mis-specified; G3b$'$ at round-off & follow-ups declared and passed\\
C07 & restart gate $10^{-6}$ failed at $6\times10^{-6}$ & chunk length, not the parent\\
C09 & checkpoint adjoint wrong by $2$--$15\,\%$ & upstream issue \#788; direct adjoint used\\\bottomrule
\end{tabular}}
\end{table}
